\section{Summary of Findings}
This thesis developed Civic.ai as a locally executable RAG-based government-service information system focused on birth registration, passport and BRTA services. Its primary contribution is the preparation and organization of a source-traceable knowledge resource. The active collection contains 3,709 chunks from 174 prepared source units, with search-oriented text separated from evidence content. Combined registration documents remain in the underlying collection, but the study does not claim evaluated death-registration capability.

The implementation adapts established RAG methods \cite{lewis2020rag,nextrag}: BGE-M3 dense retrieval, BM25, weighted RRF, reranking, automatic applicability selection and local generation. It preserves the evidence supplied to the generator and records response routes and configuration information. This makes both successful answers and failures inspectable without claiming that every source-linked response is correct.

Thirty Bangla questions were answered by CivicRAG and a matched dense-only Simple RAG baseline using Qwen3:8b. All 60 requests completed. The 180 saved-answer evaluation calculations also completed, giving mixed local self-judge results: Simple RAG has higher mean faithfulness and slightly higher answer relevancy, while CivicRAG has higher custom binary context relevance. These measures demonstrate observed behaviour, not independent factual accuracy or uniform superiority.

The separate provisional retrieval comparison provides further detail. Dense retrieval leads most overall measures, while Hybrid RRF ties dense Hit@5 overall and has higher birth-registration Hit@5 on the eligible subset. These assistant-reviewed results exclude uncertain questions and use pooled relevance, not exhaustive corpus ground truth. The overnight Qwen labeling run remains exploratory because only one question is scorable under its current rules.

The needs-analysis survey supplies a different kind of evidence: 55 of 77 valid adult willingness responses (71.4\%) favoured the proposed assistant, while 42 of 76 trust responses (55.3\%) would recheck its information. This supports the design emphasis on understandable guidance and source inspection, without demonstrating adoption or population-wide social impact.

\section{Contributions to the Field}
The principal contribution is dataset development. We organized heterogeneous civic-service materials into domain-specific retrieval records, retained original-to-derivative relationships, and preserved available metadata, URLs and unresolved review flags. This resource supports systematic investigation of civic questions rather than relying on a general model's untraceable knowledge.

The second contribution is the developed system integration. Search aliases support retrieval without becoming answer facts; chunk identifiers connect vector results to source content; and selection traces make post-retrieval decisions visible. The architecture adapts established techniques to the language and procedure distinctions of the collected documents rather than claiming a new foundation model or general-purpose RAG algorithm.

The third contribution is an auditable evaluation record. Matched baseline answers, selected contexts, model settings, local judge scores and provisional retrieval labels are retained separately. Manually selected diagnostics are not confused with automatic operation, and completed computations are not confused with validated ground truth. This separation allows future work to test specific changes without losing the current comparison.

\section{Recommendations for Future Work}
The first priority is independent, evidence-backed annotation of unseen questions and relevant passages. The set should cover paraphrases, conditional scenarios, ambiguous questions and insufficient-evidence cases across the three studied domains. Additional services, including death registration, would require explicit scope definition and evaluation rather than assuming that their appearance in a combined source proves support.

Data review should address unresolved original matches, outdated instructions, conflicting conditions and long legal OCR. Reviewed changes should be versioned with the corpus and rebuilt indexes. Relevance annotation and verification against current official documents serve different purposes and should be recorded separately.

Architectural improvements should use single-factor comparisons. Selector-on versus selector-off, adjusted context budgets and dense retrieval followed by reranking are useful candidates. Stronger multilingual generators or independent evaluation models may improve particular stages, but their benefits should be measured rather than inferred from model size or a persuasive demonstration. More capable hardware and measured runtime profiling could address the recorded interactive latency.

Finally, a prospective user study could measure clarity, usefulness and the effort needed to verify responses. Such evidence would be needed to substantiate the intended social outcomes described in the preserved Abstract. Civic.ai is presented as a developed and evaluated information system with a documented dataset contribution and bounded findings, not an official decision-making service.
